\documentclass[../../main.tex]{subfiles}% 注意这里的文件路径不能用 ./main.tex ,否则用latexmk编译子文件会报错
\graphicspath{{\subfix{./image/}}{\subfix{../../image/}}} % 指定图片引用路径,后续可以直接使用图片文件名
% 如果图片存放在上级目录,则使用 ../../image/ 作为备用路径

% 例如:
% \begin{figure}[H]
% \centering
% \includegraphics[width=0.3\textwidth]{图.png}
% \caption{}
% \label{figure:图}
% \end{figure}
% 注意:上述\label{}一定要放在\caption{}之后,否则引用图片序号会只会显示??.

\begin{document}
% SOURCE p=140; printed=125; object=ch04-nets-title; region=top
\section{网}\label{section:ch04-nets}
% SOURCE p=140; printed=125; object=ch04-nets-introduction; region=upper
如前所述，序列收敛在一般拓扑空间中，不像在度量空间中那样具有核心地位。下面的例子可以说明原因。考虑 $\mathbb R$ 上全体复值函数的空间 $\mathbb C^{\mathbb R}$，赋予乘积拓扑，也就是逐点收敛拓扑，再考虑其子空间 $C(\mathbb R)$。
% SOURCE p=140; printed=125; object=ch04-continuous-functions-sequential-density; region=upper
\begin{proposition}\label{proposition:ch04-continuous-functions-sequential-density}
$C(\mathbb R)$ 中收敛序列的极限所成的集合，是 $\mathbb C^{\mathbb R}$ 的真子集；然而 $C(\mathbb R)$ 在 $\mathbb C^{\mathbb R}$ 中稠密。
\end{proposition}
% SOURCE p=140; printed=125; object=ch04-proof-continuous-functions-sequential-density; region=upper
\begin{proof}
一方面，由\cref{corollary:ch02-complex-pointwise-limits}，若 $\{f_n\}\subseteq C(\mathbb R)$ 且 $f_n\to f$ 逐点成立，则 $f$ 为 Borel 可测函数，因而这些极限构成真子集。另一方面，对 $f\in\mathbb C^{\mathbb R}$，集合
\begin{align*}
\{g\in\mathbb C^{\mathbb R}:|g(x_j)-f(x_j)|<\varepsilon,\ j=1,\ldots,n\},\\
\qquad n\in\mathbb N,\ x_1,\ldots,x_n\in\mathbb R,\ \varepsilon>0,
\end{align*}
组成 $f$ 处的邻域基，而每个这样的集合显然都包含连续函数。
\end{proof}
% SUPPLEMENT: id=v14-ch04-nets-density-details; parent=ch04-continuous-functions-sequential-density; kind=proof-detail+motivation+hypotheses; origin=independent-derivation
\begin{proof}
\textbf{真包含与有限插值的补写（重构新增）。}连续函数 Borel 可测，复值可测函数的逐点极限仍可测，因此原证明所述的每个序列极限都 Borel 可测。另一方面，\cref{proposition:ch01-vitali-construction}所用的 Vitali 集给出一个非 Lebesgue 可测、因而非 Borel 的实数子集 $V$；若其示性函数 $\chi_V$ Borel 可测，则 $V=\chi_V^{-1}(\{1\})$ 为 Borel 集，矛盾。因此确有不属于上述序列极限集合的函数。

对稠密性，删去有限测试点中的重复项，记余下的不同点为 $t_1,\ldots,t_k$。若没有测试点，常值零函数已足够；否则定义复系数多项式
\begin{align*}
P(t)=\sum_{j=1}^k f(t_j)\prod_{\substack{1\le i\le k\\i\ne j}}
\frac{t-t_i}{t_j-t_i}.
\end{align*}
不同点保证每个分母非零，空乘积取 $1$。逐点代入可得 $P(t_j)=f(t_j)$，而 $P$ 连续，所以它属于原证明的每个指定基本邻域。任意非空开集包含这样的基本邻域，故都遇到 $C(\mathbb R)$，证明稠密性。一般有限坐标基本邻域中各坐标半径可以不同，取这些正半径的最小值便得到上面的共同 $\varepsilon$ 形式。
\end{proof}
\begin{remark}
\textbf{动机与条件说明（重构新增）。}乘积开集只检查有限个函数值，所以有限插值即可满足任何一次邻域测试；一个序列的极限却仍受到可测性的限制。这说明单靠序列无法描述一般拓扑闭包，促使我们把指标集扩大为能够同时处理各种邻域条件的有向集。
\end{remark}
% END SUPPLEMENT: id=v14-ch04-nets-density-details
% SOURCE p=140; printed=125; object=ch04-net-generalization-intro; region=upper-middle
不过，序列概念有一种能适用于任意拓扑空间的推广；关键是使用比 $\mathbb N$ 更一般的指标集。精确定义如下。
% SOURCE p=140; printed=125; object=ch04-directed-set; region=middle
\begin{definition}[有向集]\label{definition:ch04-directed-set}
\textbf{有向集}是配备二元关系 $\lesssim$ 的集合 $A$，满足：
\begin{itemize}
\item 对所有 $\alpha\in A$，有 $\alpha\lesssim\alpha$；
\item 若 $\alpha\lesssim\beta$ 且 $\beta\lesssim\gamma$，则 $\alpha\lesssim\gamma$；
\item 对任意 $\alpha,\beta\in A$，存在 $\gamma\in A$，使 $\alpha\lesssim\gamma$ 且 $\beta\lesssim\gamma$。
\end{itemize}
若 $\alpha\lesssim\beta$，也写为 $\beta\gtrsim\alpha$。
\end{definition}
% SUPPLEMENT: id=v14-ch04-directed-set-remark; parent=ch04-directed-set; kind=motivation+hypotheses+convention; origin=independent-derivation
\begin{remark}
\textbf{动机、逐条件说明与非空约定（重构新增）。}关系表示“至少进行到同样后面的阶段”。自反性把阈值本身纳入尾部，传递性允许连续推进，共同上界保证任意两个尾部有交。反复应用第三条，可给任意有限多个指标找共同上界，这使有限多项最终条件可以同时成立。不要求反对称性，因为不同指标可能代表同样精细的阶段；也不要求任意两项可比较，因为两个条件可能分别控制不同方向。

原书在列举条件时未明说 $A\ne\varnothing$。以下采用网论的通常非空指标约定，这也是后续论证实际需要的约定。若允许空有向集，则空网关于簇点的全称条件自动成立，却不可能最终进入任何邻域，下面“簇点当且仅当存在收敛子网”的结论便会失效。非空性保证“选一个指标”“取共同尾部”和收敛定义中的存在量词确有意义。后面的有限子集有向集即使底集合为空也非空，因为它仍含空子集。
\end{remark}
% END SUPPLEMENT: id=v14-ch04-directed-set-remark
% SOURCE p=140; printed=125; object=ch04-net-definition; region=middle
\begin{definition}[网]\label{definition:ch04-net}
集合 $X$ 中的一个\textbf{网}，是从某个有向集 $A$ 到 $X$ 的映射 $\alpha\mapsto x_\alpha$。通常记为 $\langle x_\alpha\rangle_{\alpha\in A}$，若指标集明确，也简记为 $\langle x_\alpha\rangle$，并称这个网由 $A$ 编号。
\end{definition}
% SUPPLEMENT: id=v14-ch04-net-remark; parent=ch04-net-definition; kind=motivation+hypotheses; origin=independent-derivation
\begin{remark}
\textbf{动机与逐条件说明（重构新增）。}网是带指标及其推进关系的映射，不能只保留取值集合；不同指标可以取相同值，重复次数和尾部位置会影响收敛。定义网本身不需要 $X$ 的拓扑，只有讨论极限时才使用拓扑。以通常顺序的 $\mathbb N$ 为指标集就得到序列；一般有向集则能让“逐渐接近”同时服从不可数多个局部条件。
\end{remark}
% END SUPPLEMENT: id=v14-ch04-net-remark
% SOURCE p=140; printed=125; object=ch04-directed-examples; region=lower
\begin{examples}\label{examples:ch04-directed-sets}
下面是有向集的若干例子。
\begin{enumerate}[label=(\roman*)]
\item\label{item:ch04-directed-integers} 正整数集 $\mathbb N$，规定 $j\lesssim k$ 当且仅当 $j\le k$。
\item\label{item:ch04-directed-punctured-line} 集合 $\mathbb R\setminus\{a\}$，其中 $a\in\mathbb R$，规定 $x\lesssim y$ 当且仅当 $|x-a|\ge|y-a|$。
\item\label{item:ch04-directed-partitions} 区间 $[a,b]$ 的所有分割 $\{x_j\}_0^n$，即 $a=x_0<\cdots<x_n=b$，规定 $\{x_j\}_0^n\lesssim\{y_k\}_0^m$ 当且仅当 $\max_j(x_j-x_{j-1})\ge\max_k(y_k-y_{k-1})$。
\item\label{item:ch04-directed-neighborhoods} 拓扑空间 $X$ 中一点 $x$ 的全体邻域组成的集合 $\mathcal N$，规定 $U\lesssim V$ 当且仅当 $U\supseteq V$。此时称 $\mathcal N$ \textbf{按反包含关系定向}。
\item\label{item:ch04-directed-product} 两个有向集的 Cartesian 积 $A\times B$，规定 $(\alpha,\beta)\lesssim(\alpha',\beta')$ 当且仅当 $\alpha\lesssim\alpha'$ 且 $\beta\lesssim\beta'$。我们总是这样把 $A\times B$ 变成有向集。
\end{enumerate}
\end{examples}
% SOURCE p=140,141; printed=125,126; object=ch04-directed-examples-discussion; region=bottom-top
前三个例子出现在初等分析中：以 $\mathbb N$ 为指标的网就是序列，第二、第三个例子中的网分别出现在实变量函数极限与 Riemann 积分的定义里。第四个例子在拓扑学中有根本重要性，而第五种构造也会多次用到。
% SUPPLEMENT: id=v14-ch04-directed-examples-proof; parent=ch04-directed-examples; kind=proof+motivation+hypotheses; origin=independent-derivation
\begin{proof}
\textbf{五例有向性的补写（重构新增）。}第一例自反、传递来自通常大小关系，共同上界可取 $\max(j,k)$。第二例自反、传递来自距离的反向大小关系；给定 $x,y\ne a$，取 $z=a+\tfrac12\min(|x-a|,|y-a|)$，则 $z\ne a$，且 $x,y\lesssim z$。两个不同点可能与 $a$ 等距，所以它说明反对称性并非必要。

第三例按严格分割约定取 $a<b$；分割族非空，各网格长度的最大值为正。自反与传递由实数大小给出。对两分割，将它们的节点合并排序得到共同加细，其每个小区间位于两原分割各自的一个小区间内，故网格宽度不大于任一原宽度，是共同上界。相同宽度的不同分割可以互相比较，因此这里也不要求反对称性。退化区间 $a=b$ 若另定义唯一退化分割及网格宽度零，也可单独处理，但不属于原式的严格分割情形。

第四例中 $X$ 是 $x$ 的邻域，故族非空；反包含关系自反且传递。任意两个邻域的交仍含 $x$ 的一个开邻域，故仍为邻域，并且是反包含关系的共同上界。第五例中 $A,B$ 均非空，积也非空；自反与传递逐坐标成立，给两对指标分别取两个坐标的共同上界，即得积中的共同上界。
\end{proof}
\begin{remark}
\textbf{动机与条件说明（重构新增）。}第二例越后面的指标越接近 $a$，第三例越后面的分割网格越小，第四例越后面的邻域越精细，第五例同时推进两个要求。第三例按网格宽度比较，并不是把“后一分割必须包含前一分割的全部节点”作为定义；共同加细只用来证明共同上界存在。第四例取邻域的交正是拓扑局部性质可转成有向结构的原因。
\end{remark}
% END SUPPLEMENT: id=v14-ch04-directed-examples-proof
% SOURCE p=141; printed=126; object=ch04-eventually-frequently; region=upper
\begin{definition}[最终进入与频繁进入]\label{definition:ch04-eventually-frequently}
设 $X$ 为拓扑空间，$E\subseteq X$。若存在 $\alpha_0\in A$，使 $\alpha\gtrsim\alpha_0$ 时都有 $x_\alpha\in E$，则称网 $\langle x_\alpha\rangle_{\alpha\in A}$ \textbf{最终进入} $E$。若对每个 $\alpha\in A$，均存在 $\beta\gtrsim\alpha$ 使 $x_\beta\in E$，则称该网\textbf{频繁进入} $E$。
\end{definition}
% SOURCE p=141; printed=126; object=ch04-net-limits-cluster; region=upper
\begin{definition}[网的极限与簇点]\label{definition:ch04-net-limits-cluster}
若对 $x$ 的每个邻域 $U$，网 $\langle x_\alpha\rangle$ 都最终进入 $U$，则称 $x\in X$ 是此网的一个\textbf{极限}，或称该网\textbf{收敛于} $x$，记 $x_\alpha\to x$。若对 $x$ 的每个邻域 $U$，该网都频繁进入 $U$，则称 $x$ 为此网的一个\textbf{簇点}。
\end{definition}
% SUPPLEMENT: id=v14-ch04-net-limit-remarks; parent=ch04-net-limits-cluster; kind=proof-detail+motivation+hypotheses; origin=independent-derivation
\begin{remark}
\textbf{动机与逐条件说明（重构新增）。}“最终”表示存在一个阈值，其后全部指标均满足条件；“频繁”表示每个阈值之后还能找到满足条件的指标。后者并不声称有无穷多个不同取值，甚至在一般有向集上也不声称有无穷多个不同指标。例如单点有向集上的网可以频繁进入包含其值的集合。

最终进入 $E$ 蕴含频繁进入 $E$：给任意 $\alpha$，取它与最终阈值 $\alpha_0$ 的共同上界 $\beta$，则 $x_\beta\in E$。又，频繁进入 $E$ 恰等价于不最终进入 $E^c$，这是将存在与全称量词逐一取否定的结果。因此收敛点一定是簇点，但簇点未必为极限，例如实序列 $(-1)^n$ 的 $1$ 是簇点而非极限。

极限与簇点只需在该点的一个邻域基上测试：任意邻域都包含基成员，最终或频繁进入较小集合便能进入较大集合。极限可能不唯一，例如两点平凡拓扑空间中的任意网收敛于两个点；唯一性需要后面的 Hausdorff 条件。网的簇点与取值集合的聚点也不同：常值实网的值是簇点，但其单点取值集没有聚点。
\end{remark}
% END SUPPLEMENT: id=v14-ch04-net-limit-remarks
% SOURCE p=141; printed=126; object=ch04-nets-three-propositions-intro; region=upper-middle
下面三个命题说明，网能很好地替代序列。
% SOURCE p=141; printed=126; object=ch04-nets-closure; number=4.18; region=middle
\begin{proposition}\label{proposition:ch04-nets-closure}
设 $X$ 为拓扑空间，$E\subseteq X$、$x\in X$。则 $x$ 是 $E$ 的聚点，当且仅当存在 $E\setminus\{x\}$ 中的网收敛于 $x$；并且 $x\in\overline E$ 当且仅当存在 $E$ 中的网收敛于 $x$。
\end{proposition}
% SOURCE p=141; printed=126; object=ch04-proof-nets-closure; region=middle
\begin{proof}
若 $x$ 是 $E$ 的聚点，令 $\mathcal N$ 为 $x$ 的邻域族，按反包含关系定向。对每个 $U\in\mathcal N$，选 $x_U\in(U\setminus\{x\})\cap E$，则 $x_U\to x$。反之，若 $x_\alpha\in E\setminus\{x\}$ 且 $x_\alpha\to x$，则 $x$ 的每个去心邻域都包含某个 $x_\alpha$，所以 $x$ 是 $E$ 的聚点。同样，若 $x_\alpha\to x$ 且 $x_\alpha\in E$，则 $x\in\overline E$；反向由\cref{proposition:ch04-closure-accumulation}得到。
\end{proof}
% SUPPLEMENT: id=v14-ch04-nets-closure-details; parent=ch04-nets-closure; kind=proof-detail+motivation+hypotheses; origin=independent-derivation
\begin{proof}
\textbf{尾部与闭包反向的补写（重构新增）。}在原构造中，给定邻域 $V$，凡 $U\gtrsim V$ 都有 $U\subseteq V$，所以 $x_U\in U\subseteq V$。故选阈值 $V$ 就验证了最终进入，而非仅仅频繁进入。反向若网收敛，对任一邻域 $V$ 存在阈值 $\alpha_0$，由自反性 $x_{\alpha_0}\in V$；其又在 $E\setminus\{x\}$ 中，因此每个去心邻域确实遇到 $E$。

对闭包反向，若 $x\in\overline E$，每个邻域 $U$ 都遇到 $E$，直接逐个选 $x_U\in U\cap E$ 即给同样的收敛网。也可沿原证明的引用：$x\in E$ 时取常值网，否则 $x$ 是聚点，应用前一结论。若 $E=\varnothing$，其闭包不含任何点，也没有非空指标集上的 $E$ 值网，两侧同时为假。逐邻域选点使用本书已采用的选择公理。
\end{proof}
\begin{remark}
\textbf{动机与逐条件说明（重构新增）。}指标现在遍历全部邻域，故无须第一可数性；这正修复了序列闭包刻画在一般空间中的缺口。聚点部分必须排除网值等于 $x$，否则常值网会使每个属于 $E$ 的孤立点也通过测试。闭包部分则应允许该值。这里没有要求网值互异，也不需要 Hausdorff 或 $T_1$。
\end{remark}
% END SUPPLEMENT: id=v14-ch04-nets-closure-details
% SOURCE p=141; printed=126; object=ch04-net-continuity; number=4.19; region=middle-lower
\begin{proposition}\label{proposition:ch04-net-continuity}
设 $X,Y$ 是拓扑空间，$f:X\to Y$。则 $f$ 在 $x\in X$ 连续，当且仅当对每个收敛于 $x$ 的网 $\langle x_\alpha\rangle$，网 $\langle f(x_\alpha)\rangle$ 都收敛于 $f(x)$。
\end{proposition}
% SOURCE p=141; printed=126; object=ch04-proof-net-continuity; region=middle-lower
\begin{proof}
若 $f$ 在 $x$ 连续，且 $V$ 是 $f(x)$ 的邻域，则 $f^{-1}(V)$ 是 $x$ 的邻域。因此，若 $x_\alpha\to x$，则此网最终进入 $f^{-1}(V)$，所以像网最终进入 $V$，即 $f(x_\alpha)\to f(x)$。

反之，若 $f$ 在 $x$ 不连续，则存在 $f(x)$ 的邻域 $V$，使 $f^{-1}(V)$ 不是 $x$ 的邻域，即 $x\notin(f^{-1}(V))^\circ$，等价地，$x\in\overline{f^{-1}(V^c)}$。由\cref{proposition:ch04-nets-closure}，存在 $f^{-1}(V^c)$ 中的网 $\langle x_\alpha\rangle$ 收敛于 $x$。然而 $f(x_\alpha)\notin V$，所以 $f(x_\alpha)\not\to f(x)$。
\end{proof}
% SUPPLEMENT: id=v14-ch04-net-continuity-details; parent=ch04-net-continuity; kind=proof-detail+motivation+hypotheses; origin=independent-derivation
\begin{remark}
\textbf{动机、条件与反向细节（重构新增）。}连续性把每个邻域条件拉回为邻域条件，所以必须保持所有网极限，而不仅是某个选定指标集的网。反向使用恒等式 $X\setminus S^\circ=\overline{X\setminus S}$，再取 $S=f^{-1}(V)$，得到原证明中的闭包。构造网的所有像值都在 $V^c$，所以它不可能最终进入 $V$：若有阈值，阈值自身已经给出矛盾。无需 $V$ 开，因为连续性与网收敛都是按邻域定义。此定理在固定点表述，不要求 $f$ 在其他点连续；也无需任一空间第一可数或 Hausdorff。
\end{remark}
% END SUPPLEMENT: id=v14-ch04-net-continuity-details
% SOURCE p=141; printed=126; object=ch04-subnet-definition; region=lower
\begin{definition}[子网]\label{definition:ch04-subnet}
网 $\langle x_\alpha\rangle_{\alpha\in A}$ 的一个\textbf{子网}，是网 $\langle y_\beta\rangle_{\beta\in B}$ 连同一个映射 $\beta\mapsto\alpha_\beta$，从 $B$ 到 $A$，满足：
\begin{itemize}
\item 对每个 $\alpha_0\in A$，存在 $\beta_0\in B$，使 $\beta\gtrsim\beta_0$ 时 $\alpha_\beta\gtrsim\alpha_0$；
\item $y_\beta=x_{\alpha_\beta}$。
\end{itemize}
\end{definition}
% SOURCE p=141; printed=126; object=ch04-subnet-limit; region=lower
\begin{proposition}\label{proposition:ch04-subnet-limit}
显然，若 $\langle x_\alpha\rangle$ 收敛于点 $x$，则每个子网 $\langle x_{\alpha_\beta}\rangle$ 也收敛于 $x$。
\end{proposition}
% SOURCE p=141; printed=126; object=ch04-subnet-warning; region=lower
\begin{remark}
使用“子网”这个名称，是因为子网发挥着与子序列十分相似的作用，但不应过于拘泥于这个名称，因为映射 $\beta\mapsto\alpha_\beta$ 未必单射。特别地，指标集 $B$ 的基数完全可能比 $A$ 更大，且一个序列的子网未必是子序列。
\end{remark}
% SUPPLEMENT: id=v14-ch04-subnet-details; parent=ch04-subnet-definition; kind=proof+motivation+hypotheses; origin=independent-derivation
\begin{proof}
\textbf{子网保极限的补写（重构新增）。}设 $U$ 为 $x$ 的任意邻域。原网收敛给出 $\alpha_0$，使 $\alpha\gtrsim\alpha_0$ 时 $x_\alpha\in U$。子网的第一条件给出 $\beta_0$，使 $\beta\gtrsim\beta_0$ 时 $\alpha_\beta\gtrsim\alpha_0$。于是这些 $\beta$ 上 $y_\beta=x_{\alpha_\beta}\in U$，故子网最终进入 $U$，证明收敛。
\end{proof}
\begin{remark}
\textbf{动机与逐条件说明（重构新增）。}第一条件说新指标最终会越过原指标集中的任意阈值，保证没有永远滞留于原网的早期部分；第二条件保证新网确实由原网取值形成。本书不要求指标映射单调、单射或满射。仅要求它的像是共尾集合还不够：对实序列 $x_n=1/n\to0$，若奇数 $k$ 取 $n_k=1$，偶数 $k$ 取 $n_k=k$，则这些指标的像无界，但 $x_{n_k}$ 不趋零，因而该指定指标映射不定义子网。

原书关于更大指标集的警告可以具体实现：在 $B=\mathbb N\times\mathbb R$ 上规定 $(n,s)\lesssim(m,t)$ 当且仅当 $n\le m$，这是非空有向集。映射 $(n,s)\mapsto n$ 满足子网条件，所以 $y_{(n,s)}=x_n$ 是原序列的子网；$B$ 不可数，映射也不是单射，这个网作为带指标的对象并不是一个序列。
\end{remark}
% END SUPPLEMENT: id=v14-ch04-subnet-details
% SOURCE p=141; printed=126; object=ch04-cluster-subnet; number=4.20; region=bottom
\begin{proposition}\label{proposition:ch04-cluster-subnet}
若 $\langle x_\alpha\rangle_{\alpha\in A}$ 是拓扑空间 $X$ 中的网，则 $x\in X$ 为此网的簇点，当且仅当此网有一个收敛于 $x$ 的子网。
\end{proposition}
% SOURCE p=142; printed=127; object=ch04-proof-cluster-subnet; region=top
\begin{proof}
若 $\langle y_\beta\rangle=\langle x_{\alpha_\beta}\rangle$ 是收敛于 $x$ 的子网，$U$ 是 $x$ 的邻域，选 $\beta_1\in B$，使 $\beta\gtrsim\beta_1$ 时 $y_\beta\in U$。又，给定 $\alpha\in A$，可选 $\beta_2\in B$，使 $\beta\gtrsim\beta_2$ 时 $\alpha_\beta\gtrsim\alpha$。存在 $\beta\gtrsim\beta_1,\beta_2$，于是 $\alpha_\beta\gtrsim\alpha$ 且 $x_{\alpha_\beta}=y_\beta\in U$。因此原网频繁进入 $U$，所以 $x$ 是簇点。

反之，若 $x$ 是簇点，令 $\mathcal N$ 为 $x$ 的邻域族，并在 $\mathcal N\times A$ 上规定 $(U,\alpha)\lesssim(U',\alpha')$ 当且仅当 $U\supseteq U'$ 且 $\alpha\lesssim\alpha'$。对每个 $(U,\gamma)\in\mathcal N\times A$，可选 $\alpha_{(U,\gamma)}\gtrsim\gamma$，使 $x_{\alpha_{(U,\gamma)}}\in U$。于是，若 $(U',\gamma')\gtrsim(U,\gamma)$，就有
\begin{align*}
\alpha_{(U',\gamma')}\gtrsim\gamma'\gtrsim\gamma,
\qquad x_{\alpha_{(U',\gamma')}}\in U'\subseteq U.
\end{align*}
由此可见，$\langle x_{\alpha_{(U,\gamma)}}\rangle$ 是原网的一个收敛于 $x$ 的子网。
\end{proof}
% SUPPLEMENT: id=v14-ch04-cluster-subnet-details; parent=ch04-cluster-subnet; kind=proof-detail+motivation+hypotheses; origin=independent-derivation
\begin{proof}
\textbf{两个量词要求的分开核对（重构新增）。}$\mathcal N$ 按反包含定向，$\mathcal N\times A$ 非空且有向，已经由前面的例子验证。簇点定义恰保证每对 $(U,\gamma)$ 的可选指标集合非空，故可以逐对选取。

为验证子网条件，固定原阈值 $\alpha_0\in A$，在新指标集取阈值 $(X,\alpha_0)$。它后面的 $(U,\gamma)$ 满足 $\gamma\gtrsim\alpha_0$，故所选 $\alpha_{(U,\gamma)}\gtrsim\gamma\gtrsim\alpha_0$。为验证收敛，给定 $x$ 的邻域 $V$，选任意固定 $\gamma_*\in A$，在新指标集取阈值 $(V,\gamma_*)$。其后的 $(U,\gamma)$ 满足 $U\subseteq V$，故所选网值在 $U\subseteq V$ 中。两个要求分别得到，不需要额外证明选择出的指标映射单调。
\end{proof}
\begin{remark}
\textbf{动机与条件说明（重构新增）。}新指标的第一坐标要求取值落入更细邻域，第二坐标要求在原网中走得更远；只用一个坐标可能丢掉另一个要求。因此乘积有向集自然同时控制趋近与子网性质。整个定理不需要第一可数性或 Hausdorff；若只允许子序列，一般簇点未必能用可数多个邻域测试到。非空指标约定在选 $\gamma_*$ 处被明确使用。
\end{remark}
% END SUPPLEMENT: id=v14-ch04-cluster-subnet-details
% SOURCE p=142; printed=127; object=ch04-net-exercises-title; region=upper-middle
\Needspace{7\baselineskip}
\subsection*{习题}
% SOURCE p=142; printed=127; object=ch04-ex30-cofinal; number=30; region=middle
\begin{exercise}\label{exercise:ch04-cofinal}
\begin{definition*}[共尾子集]\phantomsection\label{definition:ch04-cofinal}
若 $A$ 为有向集，子集 $B\subseteq A$ 满足：对每个 $\alpha\in A$，存在 $\beta\in B$ 使 $\beta\gtrsim\alpha$，则称 $B$ 在 $A$ 中\textbf{共尾}。
\end{definition*}
\begin{enumerate}
\item\label{item:ch04-cofinal-subnet} 若 $B$ 在 $A$ 中共尾，且 $\langle x_\alpha\rangle_{\alpha\in A}$ 是网，则包含映射 $B\to A$ 使 $\langle x_\beta\rangle_{\beta\in B}$ 成为原网的子网。
\item\label{item:ch04-cofinal-convergence} 若 $\langle x_\alpha\rangle_{\alpha\in A}$ 是拓扑空间中的网，则它收敛于 $x$，当且仅当对每个共尾子集 $B\subseteq A$，都存在共尾子集 $C\subseteq B$，使 $\langle x_\gamma\rangle_{\gamma\in C}$ 收敛于 $x$。
\end{enumerate}
\end{exercise}
% SUPPLEMENT: id=v14-ch04-ex30-solution; parent=ch04-ex30-cofinal; kind=solution+motivation+hypotheses; origin=independent-derivation
\begin{remark}
\textbf{定义的动机与逐条件说明（重构新增）。}共尾表示无论原指标推进到哪里，选定子集里总有更后的指标；这保留了趋近终端的可能性。对一般有向集，不能简单用“无界”替换共尾，因为可能存在互不可比方向。$B$ 使用从 $A$ 限制而来的关系。下题第二项中，$C$ 在 $B$ 中共尾且 $B$ 在 $A$ 中共尾，便推出 $C$ 在 $A$ 中共尾：先选 $\beta\in B$ 越过给定 $\alpha$，再选 $\gamma\in C$ 越过 $\beta$，最后用传递性。因此这两层的共尾要求相容。
\end{remark}
\begin{solution}
\textbf{补写解答（重构新增）。}
\begin{enumerate}
\item 先验证 $B$ 有向。因 $A$ 非空，取任一 $\alpha\in A$，共尾性给出 $B$ 中元素，故 $B$ 非空。自反与传递从 $A$ 继承。给定 $\beta_1,\beta_2\in B$，先在 $A$ 中取共同上界 $\alpha$，再由共尾性取 $\beta\in B$、$\beta\gtrsim\alpha$，便得到 $B$ 中共同上界。

对原阈值 $\alpha_0$，选 $\beta_0\in B$、$\beta_0\gtrsim\alpha_0$；凡 $\beta\gtrsim\beta_0$ 都有 $\beta\gtrsim\alpha_0$。所以包含映射满足子网的最终越过阈值条件，取值条件也直接成立。
\item 若原网收敛，每个共尾限制网由前项都是子网，因而收敛；可直接取 $C=B$。反之，假设原网不收敛于 $x$。存在 $x$ 的邻域 $U$，使原网不最终进入 $U$，故
\begin{align*}
B=\{\alpha\in A:x_\alpha\notin U\}
\end{align*}
在 $A$ 中共尾。由题设选共尾 $C\subseteq B$，使限制网收敛于 $x$。但它的全部值都不在 $U$ 中，而收敛要求最终进入 $U$；在相应阈值自身即发生矛盾。故原网必须收敛。
\end{enumerate}
\end{solution}
% END SUPPLEMENT: id=v14-ch04-ex30-solution
% SOURCE p=142; printed=127; object=ch04-ex31-subsequence-subnet; number=31; region=middle
\begin{exercise}\label{exercise:ch04-subsequence-subnet}
设 $\langle x_n\rangle_{n\in\mathbb N}$ 为一个序列。
\begin{enumerate}
\item\label{item:ch04-sequential-subnet-condition} 若 $k\mapsto n_k$ 是 $\mathbb N$ 到自身的映射，则 $\langle x_{n_k}\rangle_{k\in\mathbb N}$ 是 $\langle x_n\rangle$ 的子网，当且仅当 $k\to\infty$ 时 $n_k\to\infty$；它是\cref{section:ch00-set-language}中定义的子序列，当且仅当 $n_k$ 随 $k$ 严格递增。
\item\label{item:ch04-subsequence-cofinal-bijection} $\langle x_n\rangle$ 的子序列与由\cref{exercise:ch04-cofinal}中共尾子集所定义的子网之间，有一个自然的一一对应。
\end{enumerate}
\end{exercise}
% SUPPLEMENT: id=v14-ch04-ex31-solution; parent=ch04-ex31-subsequence-subnet; kind=solution+convention; origin=independent-derivation
\begin{remark}
\textbf{指标必须保留的澄清（重构新增）。}第一项针对题目给定的指标映射 $k\mapsto n_k$ 判断子网资格。第二项的一一对应也按“所选指标连同取值”理解；若只把相同取值的序列视为完全相同对象，则原句一般不成立。例如常值序列的偶数指标和奇数指标产生相同的取值序列，却对应两个不同共尾子集。下面明确构造指标意义下的一一对应。
\end{remark}
\begin{solution}
\textbf{补写解答（重构新增）。}
\begin{enumerate}
\item 两个指标集都是通常顺序的 $\mathbb N$，因此子网条件就是：对每个 $N\in\mathbb N$，存在 $K\in\mathbb N$，使 $k\ge K$ 时 $n_k\ge N$。这恰为 $n_k\to\infty$ 的定义。子序列的指定指标则按定义要求严格递增。严格递增的正整数序列满足 $n_k\ge k$，所以每个子序列都是子网；反向不成立，例如 $n_k=k+2+(-1)^k$ 趋于无穷，但不是严格递增。
\item $B\subseteq\mathbb N$ 共尾当且仅当它无限：有限集有最大元，不能越过它；若无限集有上界，它会包含于一个有限整数初段，矛盾。每个无限子集 $B$ 有唯一的严格递增枚举 $n_1<n_2<\cdots$：逐次取尚未选取部分的最小元。枚举穷尽 $B$，因为某个 $b\in B$ 之前只有有限多个正整数，不可能一直把它留在后面。反之，任何严格递增指标序列的值域是无限共尾子集，且它就是该集合的严格递增枚举。两构造互逆。

枚举映射 $k\mapsto n_k$ 是 $\mathbb N$ 到有序子集 $B$ 的保序双射，因而把限制网 $\langle x_n\rangle_{n\in B}$ 重新编号为相应的子序列。这正是所要求的自然对应。
\end{enumerate}
\end{solution}
% END SUPPLEMENT: id=v14-ch04-ex31-solution
% SOURCE p=142; printed=127; object=ch04-ex32-hausdorff-net-limits; number=32; region=middle-lower
\begin{exercise}\label{exercise:ch04-hausdorff-net-limits}
拓扑空间 $X$ 是 Hausdorff 空间，当且仅当 $X$ 中每个网至多收敛于一个点。（若 $X$ 非 Hausdorff，取没有不交邻域的不同点 $x,y$，并考虑有向集 $\mathcal N_x\times\mathcal N_y$，其中 $\mathcal N_x,\mathcal N_y$ 分别为 $x,y$ 的邻域族。）
\end{exercise}
% SUPPLEMENT: id=v14-ch04-ex32-solution; parent=ch04-ex32-hausdorff-net-limits; kind=solution; origin=independent-derivation
\begin{solution}
\textbf{补写解答（重构新增）。}若 $X$ Hausdorff，且某网同时收敛于不同点 $x,y$，选不交开邻域 $U\ni x$、$V\ni y$。收敛分别给出阈值 $\alpha_1,\alpha_2$，取共同上界 $\gamma$，则 $x_\gamma\in U\cap V=\varnothing$，矛盾。

反之，若 $X$ 非 Hausdorff，则存在不同点 $x,y$，使每对邻域 $U\in\mathcal N_x$、$V\in\mathcal N_y$ 都有非空交。两邻域族按反包含定向，并取乘积有向关系。对每对 $(U,V)$ 选 $z_{(U,V)}\in U\cap V$。给定 $x$ 的邻域 $U_0$，在指标 $(U_0,X)$ 以后都有 $U\subseteq U_0$，故相应网值在 $U_0$ 中，所以此网收敛于 $x$。同样，用阈值 $(X,V_0)$ 可证它收敛于 $y$。这给出具有两个不同极限的网，完成逆向。
\end{solution}
% END SUPPLEMENT: id=v14-ch04-ex32-solution
% SOURCE p=142; printed=127; object=ch04-ex33-tail-closures; number=33; region=lower
\begin{exercise}\label{exercise:ch04-tail-closures}
设 $\langle x_\alpha\rangle_{\alpha\in A}$ 是拓扑空间 $X$ 中的网，对每个 $\alpha\in A$ 令 $E_\alpha=\{x_\beta:\beta\gtrsim\alpha\}$。则 $x$ 是此网的簇点，当且仅当 $x\in\bigcap_{\alpha\in A}\overline{E_\alpha}$。
\end{exercise}
% SUPPLEMENT: id=v14-ch04-ex33-solution; parent=ch04-ex33-tail-closures; kind=solution; origin=independent-derivation
\begin{solution}
\textbf{补写解答（重构新增）。}按闭包的邻域刻画，$x\in\overline{E_\alpha}$ 当且仅当对 $x$ 的每个邻域 $U$，存在 $\beta\gtrsim\alpha$ 使 $x_\beta\in U$。要求它对所有 $\alpha$ 成立，就是
\begin{align*}
\text{对每个邻域 }U\ni x\text{ 及每个 }\alpha\in A,
\quad\text{存在 }\beta\gtrsim\alpha\text{ 使 }x_\beta\in U,
\end{align*}
这恰为簇点定义。这里只交换两个全称量词的顺序，没有交换全称与存在量词。这个刻画表示：簇点必须在任何尾部的闭包内，删除任意早期部分都不能把它排除。
\end{solution}
% END SUPPLEMENT: id=v14-ch04-ex33-solution
% SOURCE p=142; printed=127; object=ch04-ex34-weak-net-convergence; number=34; region=lower
\begin{exercise}\label{exercise:ch04-weak-net-convergence}
若 $X$ 具有由函数族 $\mathcal F$ 生成的弱拓扑，则 $\langle x_\alpha\rangle$ 收敛于 $x\in X$，当且仅当对每个 $f\in\mathcal F$ 都有 $f(x_\alpha)\to f(x)$。特别地，若 $X=\prod_{i\in I}X_i$，则 $x_\alpha\to x$ 当且仅当对每个 $i\in I$ 都有 $\pi_i(x_\alpha)\to\pi_i(x)$。
\end{exercise}
% SUPPLEMENT: id=v14-ch04-ex34-solution; parent=ch04-ex34-weak-net-convergence; kind=solution+hypotheses; origin=independent-derivation
\begin{solution}
\textbf{补写解答（重构新增）。}每个 $f\in\mathcal F$ 取值于其指定的拓扑空间 $Y_f$。弱拓扑使它们全部连续，所以原网收敛必推出每个像网收敛。

反之，设每个像网都收敛。给定 $x$ 的邻域 $U$，取包含 $x$ 且包含于 $U$ 的基本开集
\begin{align*}
W=\bigcap_{j=1}^n f_j^{-1}(V_j),
\end{align*}
其中 $V_j$ 在 $Y_{f_j}$ 中开。因为 $x\in W$，有 $f_j(x)\in V_j$。每个像网的收敛给出阈值 $\alpha_j$，其后 $f_j(x_\alpha)\in V_j$。有向性给出这有限多个阈值的共同上界 $\alpha_0$，故在 $\alpha_0$ 以后 $x_\alpha\in W\subseteq U$。若有限交为空交，则 $W=X$，任意指标都可作阈值。由此得到原网收敛。

乘积拓扑就是由坐标投影生成的弱拓扑，直接取 $\mathcal F=\{\pi_i:i\in I\}$ 即得特别情形。函数族或积指标族不需要可数，因为任一个基本邻域只使用有限多个测试。若生成族为空，弱拓扑为平凡拓扑，任意网收敛于其所在空间的每个点，恰与右侧空族全称条件相符。
\end{solution}
% END SUPPLEMENT: id=v14-ch04-ex34-solution
% SOURCE p=142; printed=127; object=ch04-ex35-unordered-sums; number=35; region=lower
\begin{exercise}\label{exercise:ch04-unordered-sums}
设 $X$ 为集合，$\mathcal A$ 是 $X$ 的所有有限子集组成的族，按包含关系定向。给定任意函数 $f:X\to\mathbb R$，对 $A\in\mathcal A$ 令 $z_A=\sum_{x\in A}f(x)$。则网 $\langle z_A\rangle$ 在 $\mathbb R$ 中收敛，当且仅当 $\{x:f(x)\ne0\}$ 是可数集 $\{x_n\}_{n\in\mathbb N}$，且 $\sum_{n=1}^\infty|f(x_n)|<\infty$；此时 $z_A\to\sum_{n=1}^\infty f(x_n)$。（参见\cref{proposition:ch00-nonnegative-sums}。）
\end{exercise}
% SUPPLEMENT: id=v14-ch04-ex35-solution; parent=ch04-ex35-unordered-sums; kind=solution+convention; origin=independent-derivation
\begin{remark}
\textbf{有限支集的澄清（重构新增）。}原题用 $\{x_n\}_{n\in\mathbb N}$ 及无穷级数书写可数支集情形。非零支集也可能为空或有限；这时应使用空和或有限和，不应重复枚举非零项来凑成一个无穷级数。准确结论是非零支集至多可数、绝对和有限；在可数无限情形才取无重复的 $\mathbb N$ 枚举。以下证明包括全部情况。
\end{remark}
\begin{solution}
\textbf{补写解答（重构新增）。}首先，$\mathcal A$ 含空集且非空，包含关系自反、传递，有限集的并仍有限并给出共同上界，所以它确为有向集。

假设 $z_A\to L\in\mathbb R$。取有限集 $S\subseteq X$，使每个有限 $A\supseteq S$ 满足 $|z_A-L|<1/2$。对任意与 $S$ 不交的有限集 $F$，有
\begin{align*}
\left|\sum_{x\in F}f(x)\right|=|z_{S\cup F}-z_S|<1.
\end{align*}
若 $F$ 只含 $f>0$ 的点，便得正项和小于 $1$；若只含 $f<0$ 的点，便得负项的绝对值和小于 $1$。因此对任意有限集 $H\subseteq X$，将 $H\setminus S$ 分成正负两部分，得到
\begin{align*}
\sum_{x\in H}|f(x)|\le\sum_{x\in S}|f(x)|+2=:C<\infty.
\end{align*}
对每个正整数 $n$，集合 $D_n=\{x:|f(x)|\ge1/n\}$ 必有限，否则可从中选多于 $nC$ 个不同点，违反上面的有限和界。非零支集 $D=\bigcup_nD_n$ 因而至多可数。若 $D$ 可数无限，无重复枚举为 $\{x_n\}$，则所有绝对值部分和均至多 $C$，由单调有界收敛，$\sum_n|f(x_n)|<\infty$。空或有限支集自然也有有限的绝对和。

反之，若支集有限，网在指标包含整个支集以后恒等于其有限和；空支集的网恒为零。若支集可数无限且绝对可和，令 $s=\sum_n f(x_n)$。给定 $\varepsilon>0$，选 $N$ 使 $\sum_{n>N}|f(x_n)|<\varepsilon$，并令 $S=\{x_1,\ldots,x_N\}$。对任何有限 $A\supseteq S$，
\begin{align*}
|z_A-s|\le\sum_{x_n\notin A}|f(x_n)|
\le\sum_{n>N}|f(x_n)|<\varepsilon.
\end{align*}
因此无序有限和网收敛于 $s$。这也识别了前向情形中的极限，因为实数极限唯一。任何另一无重复枚举都会给同一个网的极限，所以所得和与枚举无关；这里没有预先使用未经证明的任意重排保和结论。
\end{solution}
% END SUPPLEMENT: id=v14-ch04-ex35-solution
% SOURCE p=142; printed=127; object=ch04-ex36-ae-not-topological; number=36; region=bottom
\begin{exercise}\label{exercise:ch04-ae-not-topological}
设 $X$ 是 $[0,1]$ 上全体 Lebesgue 可测复值函数组成的集合。不存在 $X$ 上的拓扑 $\mathcal T$，使序列 $\{f_n\}$ 关于 $\mathcal T$ 收敛于 $f$，当且仅当 $f_n\to f$ 几乎处处成立。（使用\cref{corollary:ch02-lone-ae-subsequence}以及\subcref{exercise:ch04-cofinal}{item:ch04-cofinal-convergence}、\subcref{exercise:ch04-subsequence-subnet}{item:ch04-subsequence-cofinal-bijection}。）
\end{exercise}
% SUPPLEMENT: id=v14-ch04-ex36-solution; parent=ch04-ex36-ae-not-topological; kind=solution; origin=independent-derivation
\begin{solution}
\textbf{补写解答（重构新增）。}任何拓扑中的序列都有如下性质：若它的每个子序列都有一个收敛于同一个 $f$ 的子子序列，则原序列收敛于 $f$。这可由题目所引两道习题得到，也可直接证明：若原序列不收敛，某个 $f$ 的邻域 $U$ 外有无穷多个项，按递增指标取出这些项形成子序列；它的任何子子序列都不可能最终进入 $U$，矛盾。

现在构造打字机序列。对每个整数 $k\ge0$ 及 $j=0,\ldots,2^k-1$，定义
\begin{align*}
f_{2^k+j}=\chi_{[j2^{-k},(j+1)2^{-k})}
\quad\text{于 }[0,1].
\end{align*}
这些指标从 $1$ 开始连续排列。每个 $f_n$ 可测且属于 $L^1$，在第 $k$ 块中 $\|f_n\|_1=2^{-k}$，故 $f_n\to0$ 于 $L^1$。每个子序列仍在 $L^1$ 中趋零，由\cref{corollary:ch02-lone-ae-subsequence}，它都有一个几乎处处趋零的子子序列。

然而，对每个 $x\in[0,1)$，每块中恰有一个项在 $x$ 取 $1$；当 $k\ge1$ 时该块还至少有一个项取 $0$。故对每个这样的 $x$，序列 $f_n(x)$ 都不收敛，特别不可能几乎处处趋零。端点 $x=1$ 上各项为零，不改变这一结论。

若题述拓扑存在，每个子序列都有上述子子序列几乎处处趋零，按假设亦在该拓扑中趋零。由第一段的普遍拓扑性质，整个序列必须在该拓扑中趋零，再按假设就必须几乎处处趋零，与上一段矛盾。因此不存在这样的拓扑。这个障碍并非极限唯一性问题，证明从未要求该假想拓扑 Hausdorff。
\end{solution}
% END SUPPLEMENT: id=v14-ch04-ex36-solution
\end{document}
